\section{Study-Specific Advance Gates}
\label{app:gate-map}
These gates were written in the respective protocols before their calls. They are screening rules for further research, not interchangeable definitions of successful replication or statistical confirmation. The descriptions below locate the original criteria; the frozen protocol and executable analysis govern each result. Failure does not establish impossibility for other models, budgets or tasks.

\paragraph{Short five-trace pilot.} Under \protect\path{protocol_official_graph_five_trace_20261003.md}, inputs must qualify and all readers must normally terminate. Prefix F1 must exceed both append and fresh reading by at least five percentage points, with EM no lower and at least two of four graph differences positive against each control. This gate does not require economical deployment; all five source calls are nevertheless charged in reported cost. One length-terminated fresh reader prevents a pass.

\paragraph{Identity recoding.} Under \protect\path{protocol_official_graph_identity_recode_20261003.md}, inputs must qualify and every compact response must be valid. Compact F1 must exceed original and long-alias F1 by at least ten percentage points, with EM no lower and at least four of eight positive graph differences for each comparison. Compact model-token cost must not exceed original cost. The compact arm has one length termination, so a favorable aggregate point estimate does not pass the gate.

\paragraph{Thinking mode.} Under \protect\path{protocol_official_graph_reasoning_mode_20261003.md}, qualified inputs and at least 31 of 32 normal outputs are required for the off or fixed-route candidate. Relative to always-on, F1 may decrease by at most five percentage points, EM must not decrease, and complete token cost must fall by at least 30\%. A route dominated by always-off is not advanced. These are point-estimate screening tolerances, not formal noninferiority evidence. Both observed candidates fail.

\paragraph{Indexed source.} Under \protect\path{protocol_official_graph_indexed_source_20261003.md}, inputs must qualify and indexed-off must produce at least 15 of 16 normal outputs. Its F1 must exceed both original-on and original-off by ten points, with EM no lower; it must exceed compact-off by five points with at least four of eight positive graph differences. Total model tokens must not exceed original-off. The observed indexed-off comparison against original-on cannot satisfy the F1 condition, irrespective of the other missing arms. Compact-on was not this primary candidate.

\paragraph{Model bridge and native source qualification.} The bridge's \protect\path{protocol_official_graph_model_bridge_20261003.md} requires qualified inputs, at least 31 of 32 normal parseable 8B outputs, F1 at least $0.70$ and EM at least $0.50$. Even favorable missing-outcome bounds fail that gate. The separate native-source qualification protocol requires all sixteen outcomes and usage counts known and qualified, at least fifteen normal parseable outputs, and the same F1/EM thresholds. That study passes. These are source-competence gates, not evidence that a second-pass reader benefits from traces.

\paragraph{Native long five-trace intervention.} Its frozen protocol specifies six primary prefix comparisons: append, fresh, question-first, Re2, set mode and node majority. A practical signal needs mean F1 improvement of at least $0.05$, a positive lower descriptive paired interval and Holm-adjusted two-sided sign-test $p<0.05$. Ordering evidence concerns append; an accuracy candidate must satisfy the signal criteria against all five non-trace controls, have EM point values no lower, at least 95\% normally parsed outputs, known planned outcomes and qualified returned responses. A cost candidate additionally cannot exceed the cheapest qualified non-trace control's full logical-token cost. The interrupted archive cannot pass these gates. The all-six sign-test family must not be confused with the short-state study's five-control family.

\paragraph{Short native state.} Its prospective criteria are given in the main text. The frozen protocol is \protect\path{partial_state_protocol_20261004.md} in \texttt{recovery\_bridge\_20261003}. This study ended incomplete and did not pass. Seven of eighteen known prefix outputs were not normal-valid, making its threshold of 31/32 unattainable even if all fourteen missing outputs were valid. The prospective gate also required complete qualified execution, F1 gains of at least 0.05 over five controls with positive descriptive paired intervals and Holm-adjusted sign-test significance, EM point estimates no lower, and complete deployment-token cost no greater than the cheapest nonstate control. The post hoc visible-answer sensitivity does not alter these original criteria.
